\section{Pruning inside a high-rank classical ancestor}\label{sec:classical}

All assertions in this section concern the prepared configuration of \Cref{sec:patterns}. A structure is excluded when its full stabilizers cannot supply the failed-chart patterns required by \Cref{c:failed-pattern}; this is precisely the input needed by \Cref{c:pruning-schedule}. The statement concerns these prepared configurations, and includes the preceding exclusions specified at each stage.

\subsection{Standing interval and matrix conventions}

\begin{convention}[Prepared classical configuration]\label{c:classical-configuration}
Here we give the high-rank classical structure exclusions for the reverse-flat path argument. We work in one of the current almost simple quotients, with interval base again denoted by \(H\), top group \(M\), and simple socle \(S\le M\le{\rm Aut}(S)\). We have \(HS=M\): \(S\) is not in the base (the outer quotient is soluble and the base-core quotient must have a nonabelian monolith), and use the normal supplement property. (For a strict ancestor the top here still means that ancestor group in its almost simple quotient.) A bad structure has for label its \textbf{full stabilizer in \(M\)}. We consider labels containing the tested group and proper in \(M\); by irreducibility/fullness of the ambient simple groups below (of sufficiently high rank), the stabilizers we need are proper. Equivalently one can discard full-group stabilizers, checking properness when using the exclusions. Here is the same preparedness observation as for the witness flats in the path construction: any such proper full stabilizer containing a chart test label, or the chosen previously good anchor in the chart, for a structure on the full list of intended exclusions, also labels a trunk subflat within the relevant trial-rank range. Indeed for a strict ancestor the current path node has already avoided \emph{all} these final exclusions relative to it (even if the pruning order now starts over); for the current ancestor use properness. This applies also when testing or using a case condition on the structure before its exclusion is complete (such as which tensor side is retained). Each witness label \(K_i\) thus has full interval \([H,K_i]\) a prepared reverse link. Thus a normal subgroup of \(K_i\) is either retained (contained in \(H\)) or missing (supplemented by \(H\)), soluble normals are retained, and two commuting normals cannot both be missing (also valid modulo normal subgroups contained in \(H\)). These consequences of \Cref{f:prop-monolith} will be denoted by (\(*\)).

Recall the patterns supplied when pruning fails: many distinct labels in a private chart, including a line of \(q\) distinct labels with \(K_i\cap K_j\le K_l\) for all \(l\) on the line and \(i\ne j\) there. Moreover \(K_i\) covers \(K_i\cap K_j\), by the hyperplane-in-span property. We may require a bounded number of independent pattern points with strictly descending successive intersections, and in the initial round more distinct labels than the total height plus any fixed constant. We may use another such pattern in a separate chart for the same category, with \(K_i\cap K_j=H\) across the two. Any needed bounded subdivision into colors will be indicated; all patterns used together for a category can be taken homogeneous. Except in the initial round we use a common previously good anchor \(H<E<K_i\) for labels within a chart, with \([H,E]\) itself a prepared reverse link. Earlier exclusions imply \(E\) lies in none of the previously forbidden stabilizers. The prepared links have height exceeding any required absolute constants, interior horizontal branching, coatomisticity, and the bottom diamond parameters used in the interval reductions. Constants and bounded choices here are independent of the field and dimension of the ambient finite group; \(q\) denotes the \textbf{lattice} field size and will exceed all required absolute orbit/count bounds.
\end{convention}

\begin{convention}[Natural matrix groups and their small parameters]\label{c:classical-matrices}
Let \(V\) have dimension \(v\) over the finite ground field \(F\). The natural types are linear (\({\rm PSL}\)), unitary, symplectic, and orthogonal (nonsingular, using odd dimension only in odd characteristic); in isometry types write \(b\) for the defining nondegenerate bilinear or hermitian form, and include the defining quadratic form \(Q\) for orthogonal groups. Adjoints for \(b\) are denoted by \(*\). We use the standard descriptions of finite simple classical groups, their covering matrix groups (\({\rm SL},{\rm SU},{\rm Sp},\Omega\)), orders and automorphisms. Here \(v\) is above a sufficiently large absolute threshold. In the isometry cases \(M\) acts by semilinear similarities modulo scalars. In the linear case it acts on \(\mathcal A={\rm End}_F(V)\) by semilinear automorphisms and possibly anti-automorphisms (transpose with conjugation/field action); on the projective group the latter action uses also inverse. There is also the corresponding action on subspaces, order-reversing for an anti-action. Write \(+\) for the collineation/automorphism part (the entire group in the isometry cases). Isometry actions on \(\mathcal A\) respect the adjoint involution. All of this uses sufficiently high rank (so no triality or small exceptional automorphism issue).

We give a few more specifics on the standard classical-group facts being used (reference points are \citet{Taylor1992} and \citet{KleidmanLiebeck1990} for classical groups and their natural geometry, together with \Cref{f:conv-lie-inputs,f:conv-representation-inputs}). In recalling the group lists use parameter \(s=|F|\) except \(s=\sqrt{|F|}\) in unitary type. Forms and subspaces considered over the ground field use the finite-field Witt theory; ``nondegenerate'' for a subspace used in an orthogonal decomposition requires nondegeneracy of the restricted polar pairing also in characteristic two.

\begin{itemize}
\item In orthogonal type use \(\Omega\) as the kernel in the isometry group of determinant and spinor norm in odd characteristic, and of the Dickson invariant (rank of \(u-1\) modulo two) in even characteristic with nondegenerate quadratic space of even dimension. In particular ordinary isometry commutators are available inside it. These kernels are preserved by semilinear similarities (for the spinor norm on determinant one, scaling the form does not change the norm modulo squares, by the reflection formula with an even number of reflections). In the other types use determinant one in linear and unitary, and the full symplectic group. The quotients of linear similarity groups over the respective special groups are soluble; adjoining field automorphisms still gives soluble remaining quotients. For unitary and high orthogonal types the automorphisms needed in \(M\), including graph actions if present, are covered by the ground semilinear similarity description (no \(D_4\) here).
\item Recall for natural degree at least two the usual simplicity/perfection list: for linear and unitary the special group is perfect, simple projectively, outside degree two at \(s=2,3\) and the unitary degree three at \(s=2\); the analogous exceptions for symplectic are degree two at \(2,3\) and \({\rm Sp}_4(2)\). \(\Omega\) in the above convention/range of polar spaces has this high structure (perfect, simple projectively) in degree at least five. In degree three (odd characteristic) use the projective rank-one structure via the symmetric square of \({\rm SL}_2(s)\); in four the split determinant space has the two independent rank-one tensor actions (projectively), and nonsplit uses the natural module and its field-\(s\) twist tensoring for \({\rm SL}_2(s^2)\). Thus the split \(\Omega_4^+\) is perfect with two projective simple factors for \(s\ge4\), and \(\Omega_4^-\) is perfect, simple projectively. The plane and soluble cases needed will also be noted below.
\item The natural actions of the special groups in these assertions, apart from orthogonal planes, are absolutely irreducible, so their matrix lifts generate the full natural matrix algebra. Equivalently one can use the standard irreducible defining-characteristic vector modules and Steinberg restriction (restricted highest weight at the natural parameter; using \(B\) only for odd characteristic); in degree four orthogonal type use the rank-one tensor descriptions, and \({\rm SU}_2(s)\) gives the split \({\rm SL}_2(s)\) module over the extended field. For this last description one can write the defining hermitian form with matrix a suitable scalar times the standard nondegenerate alternating matrix in degree two (choose the scalar to make it hermitian and use equivalence of nondegenerate hermitian forms); the special group then includes \({\rm SL}_2(s)\), of the full special order. When using a derived replacement at \({\rm Sp}_4(2)\) we check it separately. The size comparisons below use only ordinary classical-group order formulas.
\end{itemize}
\end{convention}

The exclusion order is as follows, within the prepared configuration and
under the full hypotheses of the cited statements.  Initial subspace
exclusions precede direct-system pruning
(\Cref{c:linear-subspaces,c:isometry-subspaces,c:direct-pruning}).
In linear type, fields precede ordered factors
(\Cref{c:linear-fields,c:linear-factors}).  In isometry types, partial
ordered-factor pruning leaves the two stated retained-plane cases;
field transfer and prime-degree fields precede paired idempotents and
completion of the ordered-factor exclusions
(\Cref{c:partial-isometry-factors,c:field-transfer,c:isometry-fields,c:remaining-factors}).
Full spanning and elementary normalizers then precede permuted tensor
systems and the next actual simple socle
(\Cref{c:full-spanning,c:elementary-normalizers,c:tensor-pruning,c:classical-conclusion}).

\begin{lemma}[Properness of the stabilizers used below]\label{c:proper-stabilizers}
For sufficiently large ambient natural degree, each nontrivial structure used in the exclusions below has proper full stabilizer. The same properness check applies when an earlier exclusion is invoked for a structure found from the anchor.
\end{lemma}

\begin{proof}
For properness throughout, the high-rank covering matrix group of \(S\) is perfect and absolutely irreducible. It cannot preserve the initial subspace, signed-flag or individual-projector structures below (its own action is linear), nor even an adjoint-exchanged complementary pair of projectors up to swap in isometry type (perfection). It cannot normalize a proper direct-sum system below: irreducibility would force a uniform split by transitivity, then apply the direct-split order comparison below on \(V\) itself. Similarly use the tensor order comparison for ordered tensor factors or tensor systems; normalizing a field would force centralizing it by perfection, and a nontrivial elementary abelian subgroup inside \(S\) cannot be normal there. Thus when an earlier exclusion is invoked for a structure found using an anchor, its full stabilizer is indeed proper.
\end{proof}

\subsection{Initial subspace exclusions}

\begin{proposition}[Linear subspaces, signed flags and idempotents]\label{c:linear-subspaces}
In linear type, the initial round excludes stabilizers of proper nonzero subspaces when there are no anti-actions. In the presence of anti-actions it excludes signed flags \(0<I\le J<V\) and individual proper nontrivial idempotents. These exclusions require no prior anchor and make the collineation part of every subsequent good group irreducible.
\end{proposition}

\begin{proof}
First suppose linear type, without anti-actions. Exclude subspace stabilizers of nonzero proper \(U\). In each such label the shears from \(V/U\) to \(U\) are in the normal unipotent radical, hence in \(H\). Any two subspaces in the pattern must therefore be comparable. Stabilizers of a descending list along that chain have strictly descending intersections (special linear maps inside the smallest previous space can move a new proper subspace). The total interval height bounds their number, contradiction.

If \(M\) in linear type contains anti-actions, so does \(H\); fix one, \(a\in H\). Exclude the signed flag stabilizers for \(0<I\le J<V\) (automorphisms fix the two terms, anti-actions exchange them). For labels containing \(H\), \(J=aI\). Again the normal unipotent radical of the corresponding parabolic in \(S\) (preserved also by signed anti-actions) contains the shears from \(V/I\) to \(I\), and is in \(H\). Thus the \(I\)'s form a chain, the \(J\)'s the reversed chain above it. Intersections are strict taking descending \(I\)'s, just as before.

Also exclude stabilizers of individual nontrivial proper idempotents \(e\) (in the linear case with anti-actions). The two supported special linear groups on the complementary spaces are individually normal, commute projectively, and the quotient over their product is soluble. Thus exactly one is retained by (\(*\)); designate the retained space \(U\), with complement determined by \(a\) (anti-action on an idempotent sends its image to the subspace-dual transform of its kernel). In a pattern with \(\dim U\ge2\), the supported group acts irreducibly nontrivially there and fixes the complement pointwise. Any other retained space, being \(H^+\)-invariant, projects to zero along the complement or contains \(U\) (use differences under the supported group). By mutual irreducibility and distinctness, all retained spaces are pairwise in one another's complements. Their projectors are thus orthogonal idempotents in the ring sense. Intersecting the stabilizers strictly decreases until possibly the last one, by taking special linear maps on the unsplit remainder. Again this is height-bounded.

\textit{The retained line case.} For the color with retained dimension 1, take a cross-chart pair. There is a subspace \(B_0\) of codimension at most 4 in both designated complementary hyperplanes (that is, contained in both and of codimension at most 4 in \(V\)), disjoint from the span of both image lines; take a complement \(C_0\) containing both lines. The supported \({\rm SL}(B_0)\), fixing \(C_0\), belongs to \(K_i\cap K_j=H\). Thus every designated line in these patterns lies in \(C_0\) and every corresponding hyperplane contains \(B_0\), by large-dimension irreducibility. Both the total line span and the total span of dual functional lines have bounded dimension. Now a boundedly long list of independent pattern points already contradicts strict descent of intersections. Indeed descent must be strict also on intersection with \(S\), since \(HS=M\) and all groups contain \(H\). Track the growing two spans along the list; they each increase only boundedly often. Between increases, invariance is already imposed on the current spans, and further constraints are linear conditions on their endomorphism spaces (restrict matrices to the first and their transposes to the second). Each strict descent shrinks the solution subspace, by testing a witness matrix above the descent in \(S\). This bounds the length absolutely.

These exclusions imply irreducibility of the collineation part: if it preserves \(U\), then with anti-actions also \(aU\); the full group in question preserves these two with signed exchange. Nonzero intersection gives a signed flag (intersection and sum), and zero intersection a stabilized idempotent.
\end{proof}

\begin{proposition}[Singular and nondegenerate subspaces]\label{c:isometry-subspaces}
In isometry type, the initial round excludes stabilizers of nonzero totally singular subspaces, proper nonzero nondegenerate subspaces, and nonsingular radical lines in orthogonal characteristic two. It requires no prior anchor. The resulting good groups act irreducibly by semilinear transformations.
\end{proposition}

\begin{proof}
For isometry types first exclude stabilizers of nonzero totally singular subspaces \(U\) (totally isotropic for unitary or symplectic). The standard unipotent radical is again in \(H\). Here are details of the resulting comparability. Use Witt coordinates \(U\oplus Z\oplus U^-\) with last coordinate \(\lambda\) paired dually with \(U\). The radical acts by
\[
 (u,z,\lambda)\mapsto (u+f(z)+h\lambda,\ z+g\lambda,\ \lambda)
\]
with arbitrary \(f\) and accordingly adjoint \(g\), and a form-preserving correction \(h\). Such corrections exist by representing the quadratic correction as the diagonal of a bilinear pairing (orthogonal), filling an alternating correction (symplectic), or by surjectivity of quadratic-field trace on the diagonal in the hermitian case. With \(f=g=0\) the available \(h\) include the symmetric shears in symplectic type, skew-hermitian in unitary type, alternating in orthogonal type. These unipotent groups are inside the defining special groups (the additional determinant/spinor restrictions have order prime to the characteristic in the relevant unitary or odd-characteristic orthogonal cases). In orthogonal characteristic two the Dickson rank parity is even: the rank of the displacement is \(2\operatorname{rank}g\) plus the rank induced by \(h\) from \(\ker g\) to its dual \(U/{\rm im}\,f\), and the latter is alternating by the quadratic isometry condition.

Another invariant totally singular space within \(U^\perp\), if not within \(U\), must contain \(U\) by the arbitrary \(f\). A vector in such an invariant space with \(\lambda\ne0\) outside \(U^\perp\) is impossible in unitary/symplectic type by the pure \(h\)-shears (one can have nonzero pairing with \(h\lambda\)). In orthogonal type it forces the space to contain the annihilator hyperplane of \(\lambda\) in \(U\); also \(Z=0\), else choose a shear with nonsingular \(g\lambda\) and take the difference vector. Thus the only incomparable possibility consists of two split maximal singular spaces adjacent in a hyperplane. At most two distinct maxima can so occur (through their common hyperplane there are just two; a third adjacent to both and not through it contains two vectors with nonzero pairing modulo that hyperplane). The rest are in a chain. Intersections down a descending chain of singular spaces are strict by elementary special linear transvections inside the smallest previous singular space preserving the larger flag, extended dually to isometries (also in \(S\)). This again contradicts the pattern size.

Next exclude stabilizers of proper nonzero nondegenerate subspaces in isometry types. On the two orthogonal parts there are commuting normal supported special groups with soluble quotient over their product, so designate the unique retained side. For sufficiently large retained dimension its supported derived classical group acts irreducibly there, hence distinct retained spaces are pairwise orthogonal by the same mutual argument as above. Intersections strictly decrease until possibly the last (use the large irreducible supported group on the remaining orthogonal space). For bounded retained dimension use a cross-chart pair. Their two small spaces embed in a nondegenerate subspace of bounded dimension (complete the radical of their span by dual partners); the supported large irreducible classical group on its complement lies in \(H\). All the small spaces in the patterns thus have bounded total span, and strict descent for independent pattern points is bounded by the growing-span linear-constraint argument above. Use this same small-space argument for the stabilizers of nonsingular radical lines in orthogonal characteristic two. These exclusions ensure ordinary semilinear irreducibility: an invariant degenerate space has an invariant radical, totally singular except possibly in that quadratic case, where the kernel of the quadratic form on the radical is totally singular and otherwise the radical is a nonsingular line.

All initial categories needing more than height-boundedly many labels required no previously good anchor.
\end{proof}

\subsection{Direct systems of imprimitivity}

\begin{convention}[Direct systems after the subspace exclusions]\label{c:direct-configuration}
Next exclude stabilizers of sets of complementary nontrivial projectors giving a direct decomposition (at least two spaces), orthogonal nondegenerate decompositions in the isometry cases. Use the irreducible previously good \(E^+\). It is transitive on the blocks, all of common dimension \(k\), and a block stabilizer in \(E^+\) is semilinearly irreducible on its block (transport any supposed subspace). In isometry type the blocks consist of at most two isometry classes, with independent alternating groups on each class induced by the inner stabilizer: permute isometric copies, then take commutators. Similarly there is at least the alternating permutation group in linear type. If two isometry classes occur, they have equal sizes and the class partition is preserved by the full stabilizer (this uses only finite-field form types and semilinear similarities).
\end{convention}

\begin{lemma}[Two rich permutation actions]\label{c:orbit-bound}
Suppose a finite group has two permutation actions, each preserving a partition into at most two classes and inducing at least the independent alternating groups on its classes. The number of orbits on ordered pairs of pairs, with one index from each action in each pair, is bounded absolutely. Replacing indices by equivariant fibers of absolutely bounded size preserves this conclusion.
\end{lemma}

\begin{proof}
Here and below we use an orbit bound. Suppose a group has two permutation actions, each preserving a partition into at most two classes and inducing on the points at least the independent alternating groups on its classes. Then the number of orbits on ordered pairs of pairs (one index from each action per pair) is bounded absolutely. Indeed choose a size gap among the at most four class sizes, with all small sizes bounded and every large size at least 8 with alternating order exceeding the product of the quotient-index bounds from the small factorials, class permutations and parities (choose five successive absolute thresholds rapidly increasing, or more to leave a gap). Intersect the joint action with the product of all large alternating factors. This is normal of bounded index, projecting onto each large alternating factor by the index comparison and simplicity (one can restrict for this check to the preimage of the single alternating factor on its action side, whose image there is onto by hypothesis). By the subdirect description these actions split as independent diagonals, with identifications natural up to relabeling (automorphisms of alternating groups in large degree). The orbits on the specified positions are then bounded by equality patterns and small-class choices. Bounded equivariant fibers keep the bound absolute.
\end{proof}

The orbit bound will control the number of invariant partitions, rather
than the number of points. Suppose a group has at most $C$ orbits on
$\Pi\times\Pi$, and $\Omega\to\Pi$ is an equivariant map with fibers of
size at most $d$. Over a representative of each orbit on $\Pi\times\Pi$
there are at most $d^2$ ordered pairs, so there are at most $Cd^2$ orbits
on $\Omega\times\Omega$. An invariant equivalence relation is a union
of these orbits. Consequently
\begin{equation}\label{c:partition-bound}
 \#\{\text{invariant partitions of }\Omega\}\le 2^{Cd^2}.
\end{equation}
In the applications, $\Omega$ consists of common refinement cells or
tensor factors. We must recover each original system by grouping those
actual objects before this bound can count its full stabilizer labels.

\begin{lemma}[Direct line systems]\label{c:direct-lines}
Under \Cref{c:direct-configuration}, the full stabilizers of direct line systems are excluded. The only possible distinct linear systems reduce to matchings over \(\mathbb F_3\); the corresponding odd-characteristic orthogonal case also reduces to matchings. In both cases an absolute orbit bound contradicts the pattern-line size.
\end{lemma}

\begin{proof}
First exclude line systems, \(k=1\). In linear type the determinant-one diagonals in \(S\) are normal soluble in a label, thus in \(H\). A line of another system with support \(s\) on this one has diagonal orbit size \((|F|-1)^{s-1}\) (divide by \(\gcd(v,|F|-1)\) for full support), at most \(s\) by independence. For \(|F|=2\) the monomial collineations would preserve the all-ones vector line contrary to irreducibility. Thus distinct systems force \(|F|=3\) and support two throughout (\(v\) large and use \(E^+\)-transitivity). They then pair the lines of the first system by a perfect matching, putting both the sum and difference lines on each pair, uniquely. Indeed each support pair has orbit spanning its full plane, so distinct pairs are disjoint. In unitary type use determinant-one norm-one diagonals similarly; the base \(\sqrt{|F|}+1\ge3\) allows only support one. In odd-characteristic orthogonal type use sign flips even in each norm-square class, lying in \(\Omega\), again a normal soluble group. Orbit size at least \(2^{s-3}\) forces \(s\le5\); each occurring type is large, allowing arbitrary support signs, so actually \(s\le2\). Distinct systems again give matchings, uniquely per matching since the two lines on a pair must be orthogonal: they are generated by \(x+y,x-y\) from the two first-system lines with equal nonzero norms, forcing matching within one square class and the relative identification up to sign. There are no nondegenerate lines for alternating polar form. In the matching cases fix two pattern points \(i,j\) on a pattern line. Their intersection group contains alternating permutations of matched configurations independently in each of the (at most two) isometry types, so has a subgroup preserving both with absolutely bounded number of orbits on ordered pairs of first-system indices. Every other system on the pattern line gives a matching invariant by that subgroup. This bounds the number absolutely, contradiction for large lattice \(q\).
\end{proof}

\begin{lemma}[Supported block groups and exceptional planes]\label{c:direct-block-groups}
For direct blocks of dimension \(k\ge2\), the supported special groups described below give a normal projective product \(Y\), with retained block centers. Modulo those centers its nonsoluble parts are simple powers, with at most two simple sites per block. The orthogonal-plane cases in which the supported rotation group has order at most two are either already impossible or are excluded by the graph-isometry matching argument below.
\end{lemma}

\begin{proof}
We explain the supported inner group choices for \(k\ge2\). Use products across blocks of their natural special groups, actually supported (for orthogonal spaces we can take \(\Omega\) in the sense of the determinant-one spinor kernel in odd characteristic, and the kernel of the Dickson parity in even characteristic; thus in split four-space over two we use the group with the two independent natural \({\rm SL}_2\) tensor actions); replace \({\rm Sp}_4(2)\) by its simple derived group. Denote the resulting projective product by \(Y\trianglelefteq K\), and include its normal block-center subgroup among the soluble normals retained by (\(*\)). Except for soluble cases, modulo these centers \(Y\) is a direct product of isomorphic nonabelian simple groups, one per block, or two per block in orthogonal split dimension four. We recall details to control small cases. The standard natural special modules are absolutely irreducible in these dimensions except the orthogonal planes. This follows also from the usual natural modules in defining characteristic (orthogonal \(4^-\) uses the two Frobenius-conjugate natural \({\rm SL}_2\) factors tensoring after scalar extension, \(4^+\) the external tensor; orthogonal odd dimension here is in odd characteristic). For \({\rm Sp}_4(2)'\cong A_6\), restriction stays irreducible by normality in \({\rm Sp}_4(2)\) and the orders too large for \({\rm GL}_3(2)\); the end field cannot enlarge by the order in \({\rm GL}_2(4)\) either. Soluble low exceptions include natural \({\rm SL}_2\) (also \({\rm SU}_2\) or \({\rm Sp}_2\)) at parameters 2,3, \({\rm SU}_3(2)\) (its cycle and determinant-one cube-root diagonals act absolutely irreducibly), \(\Omega_3(3)\), and \(\Omega_4^+\) at 2,3; the indicated irreducibility still holds. For an orthogonal plane \(\Omega\) is the rotation group of order \((|F|\mp1)/\gcd(2,|F|-1)\), acting semisimply without fixed space or invariant nondegenerate line when that order exceeds two.

\textit{Small orthogonal planes.} For order at most two in those planes: split over two gives a canonical nonsingular vector per block, impossible by irreducibility; split over three gives a canonical orthogonal nonsingular line refinement (the two nonsingular lines in the split plane), now excluded. The remaining cases (split over five and nonsplit over three) have individual supported signs \(z_a=-1\) on block \(a\) in \(H\). Treat patterns of these cases separately. If two systems are compatible (commuting projectors), distinctness and transitivity give a common orthogonal line refinement, impossible. Otherwise some sign of one moves blocks of the other (if they all fix them the projections split the blocks). By rank two it transposes precisely two and is identity on the rest. Those two planes must be graphs \(\pm f\) of an isometry from the sign block \(U_a\) to a plane in its orthogonal complement: their sum contains \(U_a\), they project isomorphically in both directions there, and they are orthogonal. Every other sign preserves the graph pair by commutation with \(z_a\), so the partner plane is precisely some \(U_b\) (each sign outside \(a\) acts on it as a scalar). Transitivity gives a perfect matching and a graph-isometry choice on each edge. The intersection group for two fixed systems contains alternating permutations of identically aligned matched configurations (all are isometric). Every further system on the pattern line must give a matching invariant by this group, boundedly many possibilities, and boundedly many isometry choices because the field here is bounded and only one choice per edge orbit needs specifying. The edge orbit bound again follows from bounded rank on the block indices. This contradicts large lattice \(q\).
\end{proof}

\begin{lemma}[Supported groups in a prepared anchor]\label{c:supported-anchor-groups}
For every remaining direct system, each block \(a\) has an actually supported matrix subgroup \(R_a\), projectively contained in \(E^+\), of order greater than two. It is normalized semilinearly by the block stabilizer of \(E^+\). Its block module is semisimple and has no trivial constituent. This includes the case in which the supported simple power is missing from the base.
\end{lemma}

\begin{proof}
In the remaining nonexceptional direct systems, there exist for every block \(a\) actually supported matrix subgroups \(R_a\) projectively in \(E^+\), normalized semilinearly by the block stabilizer in \(E^+\), of order greater than two acting nontrivially there. For soluble \(Y\) or retained \(Y\), this follows from the full supported choices above. If \(HY=K\) is missing, \(H\) induces the full top permutation image (on blocks and simple sites), since \(Y\) acts trivially there. If two simple-site orbits occur (two per block), at most one orbit product is missing modulo centers, and the other supplies supported groups as required. Otherwise there is a transitive orbit, and we have the exact interval over \(P\) in the \(H\)-invariant subgroups of the simple power \(Y/(\text{block centers})\), where \(P\) is the base intersection image. In fact after the centers have been quotiented this transitive normal simple power is minimal normal, hence meets the base core trivially (being missing) and maps to the monolith in the core quotient. By \Cref{f:prop-interval-cases}, either \(P\) is a nontrivial product and all overgroups here are products, or all nonbottom labels are subdirect. The latter contradicts the height: strict comparable subdirect groups give strict refinement of invariant ties, whose chain length is absolutely bounded by the top orbit bound (simple sites have fibers of size at most two over blocks). Thus \(E>H\) gives an intersection product with all projections strictly larger than the corresponding nontrivial projections of \(P\). Lifting with the retained centers supplies the required singly supported \(R_a\) (one can take the full available preimage on that block in the specified supported group, including both sites if they are permuted there). Indeed all centers lost in passing from the supported choices in \(Y\) to the simple-site images have projective images in \(H\); so independent tuples in this intersection product lift with their supported representatives projectively in \(E\). Conjugating by a semilinear collineation of \(E^+\) preserving the block preserves the supported matrix choice as well, acting on representatives which are the identity on the other blocks. In all cases the block module for \(R_a\) is semisimple with no trivial constituent, since its module socle and its fixed space are invariant under the semilinearly irreducible block stabilizer.
\end{proof}

\begin{lemma}[Compatibility of direct systems]\label{c:direct-compatibility}
Any two remaining direct systems in one homogeneous pattern have commuting projectors and hence a common direct refinement.
\end{lemma}

\begin{proof}
Any two systems in such a pattern must be compatible. If all \(R_a\)'s of one fix the other's blocks individually, each projection onto \(a\) of such a block lies within the block (use differences and absence of a trivial quotient), giving compatibility. Thus incompatibility would force movement each way by supported matrices. A supported matrix moving blocks has displacement rank at most its support-block dimension and at least the moved block dimension. Hence both dimensions are equal and equality holds. On the moved system it then transposes just one pair of blocks with square the identity, and acts identically elsewhere (from the fixed dimension). Its displacement has full rank on the support block, so an involution there is impossible in characteristic two and can only be \(-1\) in odd characteristic. But a group of order greater than two cannot have only one element inducing a nonidentity block permutation, contradiction.
\end{proof}

\begin{proposition}[Completion of direct-system pruning]\label{c:direct-pruning}
The remaining direct systems are excluded, both when their supported normal product is retained and when it is missing. The bounds are absolute, independently of the ground field and natural degree.
\end{proposition}

\begin{proof}
Separate retained and missing colors. For retained systems, any two compatible systems must refine one another in some direction: incidence degrees are constant on each block side by transitivity, and if both exceed one, each touching block extends outside, hence is fixed setwise by the other's supported group there (fixed outside points). This gives an invariant proper nonzero block intersection, nondegenerate when relevant, impossible by the retained actions above. A coarse block of large dimension cannot admit the resulting uniform finer split under its full supported special group by orders. Uniformly, with base size \(s_0=|F|\) (\(\sqrt{|F|}\) in unitary type), and \(c=1\) in linear/unitary, \(c=1/2\) in orthogonal/symplectic, the large \(m\)-dimensional special group order is at least \(s_0^{c m^2-O(m)}\), whereas the isometry (or linear) stabilizer of a corresponding uniform \(t\)-block direct split is at most \(t!s_0^{c m^2/t+O(m)}\) by standard orders, smaller uniformly for large \(m\) and \(2\le t\le m\). Bounded coarse dimensions permit only boundedly many comparable refinements. This excludes the retained patterns.

For missing systems take all points of a pattern line. Compatibility
gives their common refinement into nonzero fine cells. Fix two systems
$i,j$ and group the fine cells by their pair of block indices. Each
aggregate cell has a supported special group inside $K_i\cap K_j$.
If its dimension is large, no further system can split it: a splitting
block extending outside the cell would give a proper invariant
intersection under the supported group, whereas a block contained in
the cell would give a uniform finer split contrary to the same order
bound. A small aggregate has at most its dimension many nonzero fine
cells. Thus the number of fine cells over each occupied pair is bounded
by an absolute constant $d$.

The map from a fine cell to its block-index pair is $H$-equivariant.
Indeed $H$ stabilizes both systems and permutes their common refinement;
anti-actions also permute products of commuting projectors. The two rich
top actions and \Cref{c:orbit-bound} give an absolute bound $C$ on the
orbits of ordered pairs of occupied pairs. Each original system is
recovered by grouping the fine cells and taking their direct sums.
Distinct full stabilizer labels therefore give distinct $H$-invariant
partitions of the fine cells. By \eqref{c:partition-bound} there are at
most $2^{Cd^2}$ such partitions, contradicting the chosen pattern-line
size. This completes direct-system pruning.
\end{proof}

\subsection{Tensor geometry and order comparisons}

\begin{lemma}[Ordered factors and their forms]\label{c:tensor-forms}
Let \(A\) be a unital proper nonscalar central simple \(F\)-subfactor of \(\mathcal A=\operatorname{End}_F(V)\), normalized by the current previously good anchor. Then \(V=X\otimes_FY\), with \(A=\operatorname{End}(X)\otimes1\) and commutant \(A'=1\otimes\operatorname{End}(Y)\). Both factors have dimension at least two. If \(A=A^*\), the ambient form is the tensor of factor forms. In bilinear characteristic two the factor forms are alternating; an orthogonal quadratic form is the canonical form with that polar pairing and zero on pure tensors.
\end{lemma}

\begin{proof}
We next treat subfields and ordered tensor structures. All algebras below inside \(\mathcal A\) are unital \(F\)-algebras (``prime'' extension or field algebra in these exclusions refers to prime relative degree, not to the characteristic field). Recall that a central simple matrix subfactor \(A\) gives
\[
 V=X\otimes_F Y,\qquad A={\rm End}(X)\otimes1,\qquad A'=1\otimes{\rm End}(Y)
\]
with \(A'\) its commutant (central simple algebras over finite fields split). For \(A\) proper and nonscalar both factor dimensions are at least two. Linear normalization acts by tensor matrices up to scalars, by the matrix algebra automorphism theorem; similarly with field automorphisms. If \(A=A^*\), both factors have defining forms, with \(b\) their tensor form up to a scalar: use the restrictions of the matrix involution to the two factors. These are both hermitian for unitary type, and each symmetric/alternating bilinear otherwise. Semilinear similarities preserving the ordered tensor act by semilinear similarities on the factors. In bilinear characteristic two, an ordered factor normalized by such an irreducible \(E\) cannot have nonalternating form (the diagonal contraction has an invariant hyperplane, then tensor with the other space). With both alternating, any orthogonal quadratic form in question must be the canonical one for the polar tensor, vanishing on pure tensors. Such a quadratic form exists (impose zero on basis tensors and expand, using alternation on both factors: for a pure tensor the potentially nonzero pair terms have both indices different and cancel by interchanging which second indices are matched to the two first indices); it transforms under \(E\) like the given one, and a nonzero difference, the square of a linear functional, contradicts irreducibility. We will use these facts whenever there is a common previously good anchor.
\end{proof}

\begin{lemma}[Uniform tensor order obstruction]\label{c:tensor-order}
For sufficiently large dimension \(m\), a full supported perfect natural classical group cannot normalize a proper nonscalar tensor factor or a system of at least two nontrivial tensor factors, even allowing permutations of the factors.
\end{lemma}

\begin{proof}
Useful also is the uniform order obstruction: a full supported perfect natural classical group on a space of sufficiently large dimension \(m\) cannot normalize a proper nonscalar central simple tensor factor, or a tensor system splitting the space into at least two nontrivial slots even up to permutations. Its order (also modulo scalars) is at least \(|F|^{m^2/2-O(m)}\); the linear tensor-system order is bounded by
\[
 t!\, |F|^{\sum_j x_j^2},\qquad \prod_{j=1}^t x_j=m,\quad t\ge2,\quad x_j\ge2 ,
\]
which is uniformly smaller (\(t\le\log_2 m\), \(\sum x_j^2\le m^2/4+4\), e.g. by merging factors). Use the linear action of the covering matrix group for this comparison. This proves also properness of such stabilizers in the large ambient group. Properness for field algebras follows directly from perfection (fix them pointwise if normalized by the socle) and absolute irreducibility on the natural module.
\end{proof}

\subsection{Field and ordered-factor exclusions in linear type}

\begin{proposition}[Prime-degree field algebras in linear type]\label{c:linear-fields}
After the direct-system exclusions, the full normalizers of field algebras \(D/F\) of prime relative degree are excluded in linear type, including when the ambient group contains anti-actions.
\end{proposition}

\begin{proof}
Exclude stabilizers of field algebras \(D_i/F\) of prime degrees \(r_i\). Their soluble normal projective torus \({\bf P}(D_i^*)\cap S\) is in \(H\), of order
\[
 \frac{|F|^{r_i}-1}{|F|-1}\ \big/\ \frac{\gcd(v,|F|-1)}{\gcd(v/r_i,|F|-1)}
\]
by the norm determinant. This exceeds \(r_i\) except possibly \(|F|=3,r_i=2\), when it is at least two. If a pair of fields do not commute, each torus injects by normalization into the other's relative field Galois group (a nontrivial torus element has lifts generating its prime field extension). This forces the exceptional case on both sides. If two fields commute, their generated algebra is commutative semisimple, necessarily a field since the primitive idempotent system would otherwise contradict the direct-system exclusion at \(E\). Distinct commuting fields therefore have distinct prime degrees. In a nonexceptional line take their compositum; the intersection group contains the natural determinant-one group over it, perfect if module dimension is at least two, which forces any third prime field on the line into the compositum by centralization (use elementary matrices). If dimension one, pairwise commutation alone does the same. This contradicts distinctness. In a line of the exceptional type, distinct fields must be noncommuting. They have generators \(J_i^2=-1\) with projective class in the indicated torus (norm one), and pairwise normalization gives anticommutation. Two generate a split quaternion algebra \(M_2(F)\); their intersection group includes the large determinant-one group on its tensor complement. Every line field lies in that bounded pair algebra by centralization, impossibly many. We separate the types by color here.
\end{proof}

\begin{proposition}[Ordered factors in linear type]\label{c:linear-factors}
After \Cref{c:linear-fields}, the full stabilizers of proper nonscalar ordered central simple factors are excluded in linear type.
\end{proposition}

\begin{proof}
Now exclude ordered proper nonscalar central simple factors \(A\). The projective supported special linear groups on both tensor spaces are commuting normals in the label, even allowing anti-actions, with soluble quotient. Exactly one is retained; designate its factor algebra. Its natural projective group is simple except dimension two over two or three. Over two, the characteristic order-three subgroup would give an invariant quadratic field, impossible. Over three, the \(A_4\) has a unique Klein four subgroup, normal and spanning the factor algebra by its quaternion lifts. Treat this last color separately from the simple case. For a pattern of designated algebras the groups are normal in \(H\); intersection nontrivial gives equality (or equality of the Kleins in the exceptional case), hence equality of algebras, since their lifts span. Otherwise they commute projectively and also linearly on special-group lifts (perfection, or no scalar characters for \({\rm SL}_2(3)\)). Thus distinct designated algebras commute elementwise. For a pair on a line the new algebras lie in the centralizer matrix factor, on whose space the intersection contains the full supported special linear group. In large dimension this prohibits any proper new tensor subfactor there by orders (the whole centralizer factor then permits at most one distinct new designated algebra); in bounded dimension only boundedly many distinct commuting nonscalar central simple factors can fit there anyway (their tensor-product matrix algebra embeds there by simplicity, so the product of degrees is bounded). Contradiction.
\end{proof}

\subsection{Partial ordered-factor pruning in isometry type}

\begin{proposition}[Ordered factors apart from two retained plane types]\label{c:partial-isometry-factors}
In isometry type, after the direct-system exclusions, the full stabilizers of ordered proper nonscalar selfadjoint central simple factors are excluded except possibly when the retained support is an odd-characteristic bilinear orthogonal plane or a symplectic plane over \(\mathbb F_2\). These exceptions concern the retained side of the ordered factorization.
\end{proposition}

\begin{proof}
First try to exclude ordered proper nonscalar selfadjoint central simple factors \(A\), leaving two possible remaining types temporarily. In the stabilizer the natural projective special groups supported on \(X,Y\) are commuting normals with soluble quotient. Again use \({\rm Sp}_4(2)'\) in its small case. The supports indicated lie in \(S\). Besides perfection in the usual cases, note unitary determinant one, and in orthogonal type: in characteristic two both factors are alternating as above and the tensor action from a factor has even displacement rank; in odd characteristic \(\Omega\) on a symmetric factor lies in \(\Omega\) on the tensor (diagonalize the other symmetric factor and use orthogonal sums; scaling does not change the spinor norm on determinant one). If both factors alternating in odd characteristic, the symplectic supports also lie projectively there by perfection or absence of a 2-quotient for \({\rm SL}_2(3)\). Quotient solubility follows by the two tensor projections and the standard similarity groups. Exactly one support is retained, and we designate its factor algebra.

We exclude the labels for which this retained type is neither a bilinear orthogonal plane in odd characteristic nor a symplectic plane over two. Here are the details of the supported projective groups within the designated algebra:

\begin{itemize}
\item Usually it is a single simple group with perfect special lift and absolutely irreducible spanning natural action. In split orthogonal dimension four (odd characteristic here) there are two projective \({\rm PSL}_2\) components on symplectic slots, with their matrix algebras tensoring to the factor. They are permuted by the label, even at parameter three: in the determinant-form model the singular rank-one quadric and its two rulings give these two tensor slots up to exchange (the maximal linear rank-one spaces are precisely of the two fixed-line kinds). Indeed the endomorphisms preserving each member of a ruling individually form the matrix algebra on the other coordinate (use two independent first-coordinate lines and their sum-vector line). Nonsplit four has the standard simple projective group with absolutely irreducible tensor-twist action instead.
\item The degree-two symplectic or special unitary groups at parameter three (also the symplectic slots just described at three) give projective \(A_4\) with Klein lifts spanning the slot and special lift \({\rm SL}_2(3)\) having no ground-field scalar characters. Indeed the Klein preimage in the rank-one special group at three is quaternion of order eight with commutator \(-1\) acting as negative identity in degree two, thus absolutely irreducible there (semisimplicity for this 2-group, no linear constituent possible), and the abelianization of \({\rm SL}_2(3)\) has order three whereas the ground scalar group here has order two or eight. A degree-two unitary factor at parameter two is impossible: its characteristic order-three subgroup in the projective special has two orthogonal nondegenerate eigenlines over four, giving a forbidden nondegenerate direct decomposition on tensoring. An orthogonal 3-space over three is likewise impossible: for the form up to similarity a sum of three squares, the three norm-one lines are the axes preserved as a set even by similarities, again a direct decomposition.
\item For the special unitary 3-space over four (parameter two) the projective special \(C\) of order 72 has a unique normal Sylow 3-group \(D_9\) elementary of order nine, contained in every nontrivial normal subgroup, and its lifts span the slot. Indeed in an orthonormal basis use the cycle and \({\rm diag}(1,\omega,\omega^2)\) for a cube root \(\omega\ne1\). Their commutator generates the scalar special center; projectively they give \(D_9\). In the projective full unitary group of order 216, the Fourier matrix and \({\rm diag}(1,1,\omega)\) normalize it inducing generators of \({\rm SL}_2(3)\) (rotation and shear). Thus the whole 216 normalizes it with centralizer the nine. The special image contains the nine and induces order eight acting irreducibly there (a line stabilizer in \({\rm SL}_2(3)\) has order six). Thus the nine is minimal normal and a normal subgroup disjoint from it must centralize it, proving the assertion. Its linear special lift has no character to the scalars of order three by this action and the center commutator.
\end{itemize}

These lists follow from the same natural classical groups as for direct blocks, now with odd characteristic on symmetric factors. The individual component matrix algebras here are selfadjoint.

Color according to simple components, \(A_4\) components, or the single \(C\). The projective components lie in \(H\), which permutes at most two of them in each system. In a simple or \(A_4\) pattern they cannot swap the other's pair (no quotient of order two), so components across systems mutually normalize. A nontrivial intersection, normal in each, gives equality or common Klein, hence the same component algebra; otherwise they commute projectively and linearly on special lifts as above. Likewise for the \(C\) pattern, using a common \(D_9\) if intersection is nontrivial (it contains the normal Sylow on each side, which must then be the same). Thus all the distinct component algebras are commuting tensor slots. For two labels \(i,j\) on a line their components (at most four) have a tensor centralizer factor, itself normalized by \(E\); any new component lies there. In the pair intersection is the large supported perfect classical group on that centralizer space if of large dimension, by the tensor form observations. It must stabilize individually the at most two components for each further point of the line. The order bound prohibits a proper nontrivial central simple subfactor there being new (the whole factor itself would add only one); if of bounded dimension only boundedly many commuting new ones fit. Thus the list of distinct component algebras on the whole line is absolutely bounded (at most four from the original two labels plus those new ones). Each designated algebra consists of one or two such components tensoring and determines the stabilizer (it and its tensor commutant determine the ordered tensor). Only boundedly many labels then fit on the line, contradiction. Consequently any remaining previously good-normalized proper ordered selfadjoint factorization must have the \textbf{retained} side one of the two remaining plane types.
\end{proof}

\subsection{Transfer to normalized field algebras}

The remaining field obstruction is commutation. In linear type the
retained scalar tori nearly settled this directly. In isometry type
we instead use classical groups over the extension fields. The next
setup identifies the algebra whose center determines which transfer
argument is needed.

\begin{convention}[A pair of field witnesses]\label{c:field-pair-setup}
Let $D_i,D_j$ be selfadjoint field algebras of prime relative degree,
with distinct full normalizer labels $K_i,K_j$ on a pattern line, and
put $T=K_i\cap K_j$. The algebra $R=F[D_i,D_j]$ is selfadjoint
semisimple: its radical is normalized by the common irreducible $E$
and must annihilate the space. Its primitive central idempotents are
either just one or exactly two exchanged by $*$, since sums on
$*$-orbits are orthogonal projectors contrary to the earlier
direct-system exclusion if there is more than one orbit.
\end{convention}

When the center is a field larger than $F$, or has two exchanged
components, we will find a perfect group inside $T$ whose centralizer
forces $D_i$ and $D_j$ to commute. In the scalar-center case we will
rule out $R=\mathcal A$ and use the resulting ordered tensor factor
to return to the two retained plane types.
We first establish the form and centralizer calculations needed for
these alternatives. The two-dimensional unitary exception in the
centralizer calculation will be handled by a norm-one scalar count;
exchanged components require the paired special linear group below.

\begin{lemma}[Transferred forms]\label{c:field-transfer}
Let \(D/F\subseteq\mathcal A\) be a normalized selfadjoint field algebra, of arbitrary finite degree, and put \(\delta=*|_D\). In the prepared irreducible-anchor configuration there is a nondegenerate transferred form \(B\) satisfying
\[
 b(dx,y)=\operatorname{Tr}_{D/F}\bigl(dB(x,y)\bigr).
\]
The associated bilinear, hermitian and quadratic structures, and their transformations under the normalizer, are as described below. The perfect natural special group of the transferred form, when available, lies in \(S\).
\end{lemma}

\begin{proof}
For a normalized field (degree possibly composite), let \(\delta=*|_D\). There is a nondegenerate transferred form \(B\) given by
\[
 b(dx,y)={\rm Tr}_{D/F}(d B(x,y)).
\]
It is linear over \(D\) on the first slot and \(\delta\)-semilinear on the second. It is symmetric/skew or hermitian/skew-hermitian as appropriate (skew-hermitian may be rescaled). In particular if \(b\) is alternating and \(\delta=1\) then so is \(B\), also in characteristic two by testing all square scalars. For orthogonal characteristic two and \(\delta=1\) transfer \(Q(dx)={\rm Tr}_{D/F}(d^2 Q_D(x))\), valid since the left side is additive in \(d\) (the cross term vanishes by alternation of \(B\)) and scales by squares over \(F\), so gives a linear function of \(d^2\) over \(F\); testing the trace against these same squares gives \(D\)-quadratic homogeneity of \(Q_D\) and polar \(B\). For orthogonal characteristic two and \(\delta\ne1\) use instead
\[
 Q(x)={\rm Tr}_{D_0/F}(B(x,x)),\qquad D_0=D^\delta .
\]
Indeed the right side has polar \(b\) and transforms the same way under any previously good anchor normalizing \(D\). The difference from \(Q\) has zero polar form, hence is the square of a linear functional over the finite, therefore perfect, ground field. If nonzero, its kernel would be a proper nonzero invariant subspace for that anchor. Irreducibility therefore makes the difference zero. All form transformation assertions follow from trace uniqueness under semilinear maps normalizing \(D\). For instance for a similarity \(u\) with multiplier \(\mu\) on \(b\) and induced field power \(\tau\) on \(D\), testing at \(d^\tau\) gives \(B(ux,uy)=\mu B(x,y)^\tau\). Thus the stabilizer acts projectively by corresponding \(D\)-semilinear similarities (extra scalar kernel soluble), and the standard \(D\)-linear special group, when perfect, lies inside \(S\).
\end{proof}

\begin{lemma}[Centralizers of perfect transfer groups]\label{c:transfer-centralizer}
Suppose the field normalizer in \Cref{c:field-transfer} is nonsoluble and the total ground-field dimension is sufficiently large. There is a perfect natural classical group \(P_D\), absolutely irreducible over \(D\). Its centralizer in ground-field endomorphisms is exactly \(D\), except possibly for a unitary \(D\)-module of dimension two with \(\delta\ne1\) fixing \(F\). This exception does not occur for prime extension degree in sufficiently large total dimension.
\end{lemma}

\begin{proof}
If this normalizer is nonsoluble and total ground-field dimension large, we can use a perfect such classical group \(P_D\) on the natural \(D\)-module, absolutely irreducible over \(D\). Indeed either module dimension or extension degree/field size is large, avoiding the small parameter exceptions; in the permanently soluble form dimensions nonsolubility of the normalizer would be impossible. Its centralizer in ground-field endomorphisms is exactly \(D\), except possibly the unitary module of dimension two over \(D\) with \(\delta\ne1\) fixing \(F\). To see this, extending from \(F\) to algebraic closure splits the module into copies twisted by the different embeddings of \(D\) over \(F\). They are irreducible and pairwise inequivalent since matrix traces generate \(D\) over \(F\): symplectic type uses an \({\rm SL}_2\) block; orthogonal type can use split 4-space with tensor trace there an arbitrary product, or norm traces from the \(4^-\) tensor twist, and in dimension three the symmetric-square traces (odd characteristic). Unitary type has traces generating at least \(D_0\) via \({\rm SU}_2\), and in dimension at least three a norm-one diagonal with entries \(a,b,(ab)^{-1}\), where \(a,b,ab\ne1\), has trace differing from inverse trace, giving all of \(D\). In dimension two, \(D_0F=D\) if \(\delta\) is nonidentity already on \(F\). Here on a split nondegenerate orthogonal 4-space use the two \({\rm SL}_2(D)\) tensor actions, supported there (available also within higher dimensions of Witt index at least two); in the nonsplit dimension-four case the trace on the tensor twist ranges over norms of arbitrary \({\rm SL}_2\)-traces, and in dimension three the values \(u^2-1\) from the symmetric square generate the field in odd characteristic. For \({\rm SU}_2\) the natural action over \(D\) realizes the split rank-one traces over \(D_0\). The indicated three-entry unitary diagonal exists even when the norm-one group has order three (take \(a=b\)); the difference of its trace from inverse trace, multiplied by \(ab\), is \(-(a-1)(b-1)(ab-1)\). Thus the embedding distinction (and absolute irreducibility in each copy, as for natural classical groups above) gives scalar centralizer on each over closure whenever claimed, hence the centralizer assertion over \(F\). In particular the exception does not arise for prime degree in sufficiently large ground-field dimension (it would require degree two and module dimension two).
\end{proof}

\begin{lemma}[Paired transfer spaces]\label{c:paired-transfer}
In the high-dimensional prime-degree setting of \Cref{c:field-pair-setup}, suppose an unordered complementary pair \(e,e^*=1-e\), normalized by \(E\), commutes with \(D_i\). Its two spaces are dually totally isotropic and, in orthogonal type, totally singular. Their simultaneous field-and-pair stabilizer contains the paired group \(\operatorname{SL}_m(D_i)\). In the nonsoluble case this group is perfect in \(S\), and every ground-field endomorphism centralizing it commutes with \(D_i\).
\end{lemma}

\begin{proof}
We will also use the following paired situation (still with the prime-degree \(D_i/F\) and high total dimension). Suppose two central or commuting projections to be used below are exchanged by \(*\) (a complementary pair \(e,e^*=1-e\) normalized as a pair by \(E\)), and commute with \(D_i\). Their spaces are dually totally isotropic, and totally singular also for \(Q\) when relevant (in characteristic two each quadratic restriction with zero polar is a linear square, and the sum of the two kernels would be a proper nonzero invariant space if a restriction were nonzero; the dimensions are large). The full simultaneous stabilizer of \(D_i\) and the pair thus includes paired \({\rm SL}_m(D_i)\) on these spaces, acting by inverse adjoint on the second, where \(m\) is their \(D_i\)-dimension; if \(m=1\) the stabilizer would be soluble. For \(m\ge2\) this group is perfect inside \(S\). Every endomorphism centralizing its matrices commutes with \(D_i\). Indeed over closure the constituents are natural at embeddings \(\tau\) over \(F\) and dual at \(\tau'\delta_i\) with \(\tau'\) over \(F\) (by the transferred pairing); the field scalars act there by \(\tau,\tau'\), respectively. Natural twists are distinguished by traces, as are dual twists; a natural and dual identification could only have \(\tau=\tau'\delta_i\) by traces on an \({\rm SL}_2\) block. It cannot occur for \(m\ge3\) (diagonals \(a,b,(ab)^{-1}\) again; \(|D_i|\ge4\)). If \(m=2\), either the two identified constituents have the same field-scalar action, or \(\delta_i\ne1\) fixes \(F\), forcing prime degree two and bounded total dimension. This proves the claim.
\end{proof}

\subsection{Prime-degree field pruning in isometry type}

\begin{proposition}[Prime-degree selfadjoint fields]\label{c:isometry-fields}
After direct-system pruning and the partial ordered-factor exclusion of \Cref{c:partial-isometry-factors}, the full stabilizers of selfadjoint field algebras of prime relative degree are excluded. The pattern is colored by whether the adjoint restriction on the field is the identity.
\end{proposition}

\begin{proof}
\textit{A larger center field.} If the center of \(R\) is a field larger than \(F\), take a central prime extension \(D'\) invariant under \(*\) and \(T\). If \(K_i\le N_M(D')\), then \(P_{D_i}\) forces \(D'\subseteq D_i\) by centralization, so \(D_i=D'\) commutes with \(D_j\). Otherwise the covering property gives
\[
 T=K_i\cap N_M(D')
\]
exactly. Consider the commutative semisimple \(L_f=F[D_i,D']\). If a field, the degrees must be distinct primes with compositum degree their product, and the perfect transfer group over it is available in \(T\) (the subfields are unique, and \(T\) is nonsoluble since it contains \(E\)). This perfect group centralizes \(D_j\), forcing commutation with \(D_i\) by the centralizer assertion. In the exceptional unitary dimension two case here, the compositum degree is \(2s\) with \(s\) a large odd prime (total dimension \(4s\)). Write \(a=|F|\) and identify the compositum with \(C=\mathbb F_{a^{2s}}\). Work in the actual scalar group
\[
 U=\{z\in C^\times:z^{a^s+1}=1\},\qquad |U|=a^s+1.
\]
The subgroup \(U_S\) whose projective classes lie in \(S\) has index at most two in \(U\): these are actual determinant-one isometries, with at most the spinor or Dickson restriction still to impose. Their classes belong to \(T\). Let \(U_0\le U_S\) be the scalar preimage of the kernel of their action on \(D_j\). Since the matrices are \(F\)-linear, this action has order at most \([D_j:F]\). That prime degree divides \(4s\), so it is either \(2\) or \(s\), and
\[
 |U_0|\ge\frac{a^s+1}{2s}.
\]
The two maximal proper relative subfields of \(C/F\) are \(\mathbb F_{a^s}\) and \(\mathbb F_{a^2}\). Their intersections with \(U\) have sizes at most \(2\) and \(a+1\), respectively: the first satisfies \(z^2=1\), while on the second the oddness of \(s\) turns the norm-one condition into \(z^{a+1}=1\). For sufficiently large \(s\), the displayed lower bound exceeds \(2+(a+1)\). Hence \(U_0\) contains a generator of \(C/F\). Ground-field scalar changes do not affect conjugation on \(D_j\), so this is a count in the scalar preimage throughout. Such a generator centralizes \(D_j\), forcing all of \(C\), and in particular \(D_i\), to commute with \(D_j\).

\textit{A split compositum.} If \(L_f\) here splits, there are exactly two exchanged components, and \(D_i,D'\) must have the same prime degree \(r\) (tensor over \(F\)). The possible primitive components are indexed by relative identifications of the two degree-\(r\) fields; \(E\) acts on these by translations of a cyclic \(r\)-set (the two field automorphisms extend the same base automorphism), transitively on the occupied subset. Thus \(r=2\). Now apply the paired situation to the two components: the indicated paired special linear group also centralizes \(D'\) (on each component that field acts through \(D_i\)). Hence it lies in \(T\), centralizes \(D_j\), and gives commutation with \(D_i\).

\textit{Two exchanged central components.} If the center of \(R\) instead has two exchanged components, use their projection-pair stabilizer in place of \(N_M(D')\). It cannot contain \(K_i\) by perfection and the centralizer of \(P_{D_i}\), so its full intersection with \(K_i\) is \(T\), by covering. The paired situation applies directly, yielding the same commutation.

\textit{Scalar center.} It remains to consider center \(F\). Here \(R\ne\mathcal A\), else \(T\)'s actions on the two fields have trivial projective kernel and soluble image. Thus \(R\) gives a proper ordered selfadjoint factorization, necessarily with one of the retained remaining plane factors by the partial exclusion. If \(R\) itself is that plane matrix algebra, then for bounded field size the intersection contains the high-dimensional perfect classical group on its tensor complement, forcing every line field into bounded \(R\), contradiction. For an odd orthogonal plane factor at large field size the projective supported rotations (from \(\Omega_2\)) have order greater than four, and lie in \(H\) as a normal soluble group for the stabilizer label of \(R\) (which contains \(E\)). Normalizing both \(D_i,D_j\) (quadratic within \(R\)), some nonscalar rotation centralizes both. In the plane matrix algebra this gives their equality (the centralizer of a nonscalar matrix there has dimension two), impossible. If instead \(R'\) (the matrix commutant of \(R\)) is the remaining plane algebra, \(K_i\) cannot preserve \(R\), since then \(P_{D_i}\) would act trivially on \(R'\) by its soluble similarity projection, contrary to its field centralizer. Thus again \(T\) is the full intersection \(K_i\cap N_M(R)\). On the large \(R\)-factor space apply transfer to \(D_i\) there, selfadjoint for its factor form. Its field normalizer there must be nonsoluble (otherwise both factor projections of \(T\) soluble). We get its supported perfect transfer group in \(T\), with centralizer in \(R\) exactly \(D_i\); in characteristic two here use the alternating factor form, the total tensor preserving its given orthogonal quadratic if necessary as explained earlier. This group centralizes \(D_j\subset R\), impossible.

\textit{Completion of the line argument.} Consequently the fields on a pattern line all commute unless we already obtained a contradiction. Two distinct ones of the same prime degree would give a split commutative algebra with exactly two exchanged components as above, degree two, and the two adjoint restrictions must differ (to exchange the relative-identification idempotents). Use homogeneous color of the adjoint type (identity or not on the prime-degree field algebras) to prevent this. Thus two have distinct primes and generate a field. Its transfer group in the nonsoluble intersection forces any third into the compositum, impossible by prime degrees. There is no unitary dimension-two exception here: if the involution fixes \(F\) but not the compositum, the two prime fields have different adjoint types. This completes prime-field pruning.
\end{proof}

\subsection{Paired idempotents and the remaining ordered factors}

\begin{proposition}[Completion of the algebra exclusions]\label{c:remaining-factors}
After \Cref{c:isometry-fields}, the full stabilizers of unordered proper idempotent pairs \(e,e^*=1-e\) are excluded. The two remaining retained-plane types in \Cref{c:partial-isometry-factors} are then impossible, completing the ordered-factor exclusions.
\end{proposition}

\begin{proof}
Exclude stabilizers of unordered pairs \(e,e^*=1-e\) of proper idempotents. For two such labels, choosing idempotents \(e,f\), the generated selfadjoint algebra is again semisimple. A center field larger than \(F\) is now forbidden. If there are two central components exchanged, their projector pair commutes with each of the given pairs. Distinct commuting such pairs are impossible: they would give all four projector cells (or already be the same pair, using adjoint exchange), with forbidden orthogonal direct sums on adjoint orbits. For center \(F\), note \((e-f)^2\) is central, and the algebra is spanned over this scalar algebra by \(1,e,f,ef\), since the relation rewrites \(fe\) in terms of \(e,f,ef\) and scalars and the generators are idempotent. It is therefore \(M_2(F)\), which must have remaining plane type. An orthogonal plane has at most one such pair. For the bounded symplectic plane case, the high perfect tensor-complement group in the intersection fixes individually each further pair on the line, forcing them into the bounded plane algebra. Thus all paired stabilizers are excluded.

Finally the remaining ordered factors themselves cannot occur: an odd orthogonal plane factor has the unique pair of isotropic lines over the closure, giving a forbidden projector pair if split, or a forbidden quadratic field (the diagonal algebra descended to \(F\)) if nonsplit, selfadjoint and normalized by all the factor similarities. A symplectic plane over two gives a normalized selfadjoint quadratic field via its characteristic order-three subgroup. This completes the stated field and ordered-factor exclusions.
\end{proof}

\subsection{Full spanning and elementary normalizers}

\begin{lemma}[Full spanning]\label{c:full-spanning}
After the subspace, direct-system, field, ordered-factor and paired-idempotent exclusions, every nontrivial projective subgroup \(P\le S\) normalized by a previously good anchor satisfies \(F[\text{lifts of }P]=\mathcal A\).
\end{lemma}

\begin{proof}
After the preceding exclusions we have the following full spanning property: any nontrivial projective subgroup \(P\le S\) normalized by a previously good anchor has
\[
 F[\text{lifts of }P]=\mathcal A .
\]
Indeed the radical vanishes by anchor irreducibility, and in isometry type the algebra is selfadjoint. Multiple center components are now impossible by the direct-system exclusions (and the exchanged pair exclusion for adjoints). Center-field enlargement and a proper nonscalar central simple factor are also excluded. Semilinear anti-actions in linear type respect the algebra as well (their group action on lifts uses inverse, harmless here).
\end{proof}

\begin{proposition}[Elementary abelian normalizers]\label{c:elementary-normalizers}
Under the full-spanning conclusion of \Cref{c:full-spanning}, full normalizers of nontrivial elementary abelian subgroups of \(S\) are excluded.
\end{proposition}

\begin{proof}
Exclude full normalizers of nontrivial elementary abelian subgroups \(P\le S\). In a pattern each such subgroup is contained in \(H\le E\) by (\(*\)). Its scalar commutator pairing from linear lifts is nondegenerate by full spanning. Thus it identifies \(P\) with its ground-field character dual (nondegeneracy already forces the exponent prime to divide \(|F^\times|\)). Lifts are linearly independent by their distinct simultaneous conjugation characters, and their linear span is the generated algebra. Hence \(|P|=v^2\) and its projective linear centralizer is itself (match the scalar conjugation character using the pairing, then use full spanning). There cannot be two distinct ones here normal in \(E\): a nontrivial intersection would itself satisfy the same properties, and a trivial intersection would give projective commutation. This excludes the pattern.
\end{proof}

\subsection{Tensor systems and missing normal powers}

\begin{lemma}[Reduction to one simple group per tensor slot]\label{c:tensor-simple-slots}
After the preceding exclusions, a tensor system has transitive equal-dimensional slots and the appropriate tensor forms. Its stabilizer induces at least the alternating group on the slots. To exclude all tensor systems, it suffices to exclude those having a normal supported projective product \(Y=B^t\le K\cap S\), with one nonabelian simple group per slot whose lifts span the slot algebra.
\end{lemma}

\begin{proof}
Now exclude stabilizers of tensor systems of at least two nontrivial matrix factors tensoring to \(\mathcal A\), allowing permutations of slots, with factors individually selfadjoint in isometry type. For systems under consideration the anchor acts transitively, else an orbit product contradicts ordered-factor exclusion. Thus we have \(t\) slots of common dimension \(k\). The tensor forms in isometry type can be aligned to identical forms on the factors up to scalars by transitivity/similarity (finite-field form classes). In bilinear characteristic two all are alternating (by transitivity and alternation on \(V\)); again any orthogonal quadratic here is the canonical one zero on simple tensors by irreducibility. Consequently the stabilizer label \(K\) induces at least \(A_t\) on the slots by taking commutators of actual tensor-slot permutations inside the full linear or isometry group.

We can reduce to the case of a normal supported projective product \(Y=B^t\) inside \(K\cap S\), one nonabelian simple group per slot (embedded projectively), whose linear lifts span each slot matrix algebra. Use the natural special groups as above, including \({\rm Sp}_4(2)'\). Here is the treatment of other possibilities. Small degree-two parameter cases at two and three in the linear, symplectic and unitary families give canonical nontrivial elementary abelian projective groups per slot (order three or the Klein, characteristic in the projective special). Their product is elementary abelian projectively inside \(S\), now forbidden; the same applies for \({\rm PSU}_3(2)\) using \(D_9\). Containment uses the same special supports as for ordered factors. \(\Omega_3(3)\) again gives direct orthogonal line cells from the axes. Odd-characteristic orthogonal planes of split type give pairs of singular lines per slot, whose tensor cells paired with their duals give a forbidden nondegenerate direct decomposition. In nonsplit type use the canonical quadratic fields per slot; their tensor algebra has \(2^{t-1}\) components, each central projector selfadjoint (conjugation on all the fields simultaneously), again forbidden. Split orthogonal four-space factors canonically refine into pairs of symplectic plane tensor slots as above, giving a normalized tensor system of the latter instead. Thus it suffices to exclude the single-simple-per-slot systems (first treating just these labels; the refinement observation then handles the four-space case).
\end{proof}

\begin{lemma}[Retained and missing tensor products]\label{c:tensor-products}
Retained products from \Cref{c:tensor-simple-slots} give at most one system in a pattern. For a homogeneous missing pattern, every \(N_i=E\cap Y_i\) is a nontrivial slotwise product \(\prod_aR_{i,a}\), strictly above the base intersection.
\end{lemma}

\begin{proof}
For a pattern with \(Y\) retained, these products are minimal normal subgroups of \(E\) by slot transitivity. Two must be equal (commutation impossible by perfection and spanning), which recovers the same system by the lifts of the simple factors. For missing patterns \(HY=K\), \(H\) induces the full top action, transitive and containing the natural alternating group. Exactly as for missing direct simple-site products, the interval over \(H\cap Y\) in \(Y\) of invariant subgroups must be the nontrivial product case by the reductions and subdirect tie-chain bound. In particular on a pattern line, for each system \(i\),
\[
 N_i=E\cap Y_i=\prod_a R_{i,a}\trianglelefteq E,
\]
a slotwise projective product strictly above the base intersection.
\end{proof}

\begin{lemma}[Common refinement of missing tensor systems]\label{c:tensor-refinement}
A finite pattern line of the missing systems in \Cref{c:tensor-products} has a common tensor refinement formed from the lift algebras of tuple intersections \(\bigcap_iR_{i,a_i}\). After scalar factors are removed, those algebras tensor to \(\mathcal A\) and group to recover every original slot. They are permuted transitively by \(E\). In bilinear characteristic two their forms are alternating; if an orthogonal quadratic form is present, it is canonical also for every grouping into two nonempty batches.
\end{lemma}

\begin{proof}
Write bars for the quotient by $\operatorname{core}_E H$, and let $W$
be the nonabelian simple-power monolith of $\bar E$. Each $\bar N_i$
contains $W$. The simple factors of $W$ partition among the
$\bar R_{i,a}$: their commutators with $W$ generate $W$ by perfection,
lie in the respective intersections with $W$, and are normal there and
mutually commuting.

We need these assignments to give nontrivial intersections in the actual
group, not only compatible images in the core quotient. Put
$N=\bigcap_iN_i$. Since the $N_i$ are normal in $E$, an iterated
commutator of them lies in $N$; its image contains $W$ by perfection.
Thus $\bar N\supseteq W$. Fix a simple factor $C$ of $W$, and let $a_i$
be its assigned slot in system $i$. Because $R_{i,a_i}\trianglelefteq N_i$
and $N\le N_i$, we have
\[
 [N,R_{i,a_i}]\le N\cap R_{i,a_i}.
\]
The image of this commutator contains $C$. Each $N\cap R_{i,a_i}$ is
normal in $N$, so a second iterated commutator shows that
\[
 M_\gamma:=\bigcap_iR_{i,a_i},\qquad \gamma=(a_i),
\]
has image containing $C$. In particular it is nontrivial.

Now take the intersections $M_\gamma$ for all slot tuples, allowing
trivial intersections. They are permuted by $E$ and generate a
nontrivial $E$-normalized subgroup of $S$, whose lifts span
$\mathcal A$ by \Cref{c:full-spanning}. For distinct tuples the lifts
commute linearly, because they lie in distinct tensor slots of at least
one original system. Their lift algebras therefore commute and jointly
generate $\mathcal A$. The center of each algebra centralizes all of
$\mathcal A$, hence is scalar. Its radical, multiplied by the commuting
remaining algebras, would give a nilpotent ideal of $\mathcal A$, hence
vanishes. The non-scalar lift algebras are thus central simple; their
tensor product maps isomorphically onto $\mathcal A$.

This also recovers the original embedded factors. Group the tuple
algebras by their assigned slot in any one system. Each grouped algebra
is contained in that slot, and the product of all grouped algebras is
still $\mathcal A$. Comparing dimensions in the original tensor
decomposition forces equality in every slot. We have obtained a common
tensor refinement, not merely abstract isomorphic factors.

The fine factors are selfadjoint when appropriate and are permuted
transitively by $E$, by the ordered-factor exclusion. Thus the same form
and quadratic observations apply. In particular in bilinear
characteristic two the fine-slot forms are all alternating, by
transitivity and alternation of the total form. Grouping them into two
nonempty batches still gives alternating forms on both. The orthogonal
quadratic, if present, is zero also on pure tensors for every such
two-batch grouping: its canonical quadratic with the given polar form
exists by alternation and is zero on fine pure tensors too. The two
quadratics therefore agree on a basis and, having the same polar form,
agree everywhere.
\end{proof}

\begin{proposition}[Completion of tensor-system pruning]\label{c:tensor-pruning}
The missing tensor systems of \Cref{c:tensor-products} are excluded by an absolute bound on the possible invariant partitions of their common fine slots. Hence all tensor systems are excluded.
\end{proposition}

\begin{proof}
Fix two systems $i,j$ and send each nontrivial fine factor to its pair
of original slot indices $(a_i,a_j)$. Grouping factors with the same
pair gives an aggregate matrix factor. On a large aggregate space,
$K_i\cap K_j$ contains the full supported perfect classical group.
No further system on the line can split this aggregate: its fine
factors would give a tensor system normalized by that group, contrary
to \Cref{c:tensor-order}. Thus a large aggregate contains only one
nontrivial fine factor. On a small aggregate of degree $h$, there are
at most $\log_2h$ such factors, because their degrees multiply to $h$
and each is at least two. The fibers of the displayed map are therefore
bounded absolutely.

The map is $H$-equivariant, since $H\le E$ stabilizes both original
systems and permutes all tuple intersections; discarding scalar factors
preserves this action. By \Cref{c:orbit-bound}, the two rich top actions
give an absolute bound on ordered-pair orbits of occupied pairs, hence
on ordered-pair orbits of fine factors. Each original system is recovered
by grouping those embedded factors, as proved in
\Cref{c:tensor-refinement}. Distinct stabilizer labels therefore give
distinct invariant partitions. The bound \eqref{c:partition-bound}
contradicts the pattern-line size and completes tensor-system exclusion.
\end{proof}

\begin{theorem}[The next classical-ancestor socle]\label{c:classical-conclusion}
Let \(K\) be a tested descendant with a prepared interval, good for all the exclusions in this section, or containing the final good anchor. There is an actual nonabelian simple subgroup \(B\trianglelefteq K\) inside the ancestor socle \(S\), with full natural spanning, such that
\[
 HB=K,\qquad C_K(B)\le\operatorname{core}_K H,
\]
and \(C_K(B)\) is soluble. The quotient by this centralizer is almost simple with socle \(B\), and the exact subgroup interval is preserved.
\end{theorem}

\begin{proof}
Finally consider a tested descendant \(K\) with prepared interval, good for these exclusions (or containing the final good anchor). It has \(K\cap S\ne1\) since \(K\) is nonsoluble. Take a minimal \(K\)-normal subgroup there. It cannot be elementary abelian, and if a power of nonabelian simples, their lift algebras commute linearly by perfection, and jointly span fully. As in the refinement argument they are central simple tensor factors, individually selfadjoint when relevant; multiple ones would be a forbidden tensor system. Thus there is an actual simple \(B\trianglelefteq K\) in \(S\), absolutely irreducible projectively. Its projective linear centralizer is trivial by perfection and spanning, so \(C_K(B)\) injects in the soluble outer quotient and lies in \({\rm core}_K H\) by (\(*\)). Also \(B\) is not retained (else base-core quotient soluble), hence \(HB=K\). Quotienting its centralizer gives the required next almost simple group with socle \(B\), keeping the exact interval.
\end{proof}
